\section{Discussion}
\label{sec:discussion}

\subsection{Summary of findings}

Our simulations indicate that a revenue-neutral replacement of council tax with a 0.77 per cent flat land value tax would leave 68 per cent of households better off, reduce absolute poverty after housing costs by 0.70 percentage points, and shift the tax burden substantially up the wealth distribution, while leaving before-housing-cost poverty and the Gini coefficient of equivalised household net income essentially unchanged. The dissociation between the wealth-decile and income-Gini results is the paper's main analytical point, and it generalises: distributional statistics computed over the income distribution can be nearly uninformative about a reform whose incidence is organised by wealth. An LVT is a tax on a stock that is highly concentrated --- we estimate that the top wealth decile holds 47.4 per cent of all land, and the wealth Gini is 0.70 against 0.36 for income --- assessed on households whose position in the income distribution is only weakly related to that stock. A reader who evaluated this reform only on the Gini and the BHC poverty rate, the two statistics that dominate UK distributional debate, would conclude that it does almost nothing. Table~\ref{tab:wealth_impact} shows that it does a great deal.

We are correspondingly cautious about the income-decile profile. Once the top two deciles are set aside, the average effects in deciles 2--8 are a few hundred pounds a year on incomes of \pounds 31{,}000--\pounds 69{,}000, they are not monotonic in income, and they rest on imputed land values whose decile means are visibly noisy. We do not think the data support claims about which particular middle decile gains; they support the claim that the reform is roughly neutral across the bottom eight deciles and costly to the top two.

The efficiency case, which our static framework does not quantify, points in the same direction as the equity case. Whereas council tax is levied on the value of dwellings, including improvements, an LVT imposes no charge on constructing, extending or upgrading property, and its base is immobile \citep{george1879, oates1997}.

\subsection{Limitations}

Several caveats qualify the results.

\paragraph{Static incidence.} We assign the full burden to current landowners, with no behavioural response and no capitalisation in the central results. Standard theory predicts that an anticipated recurrent land tax capitalises into lower land prices \citep{feldstein1977}, so part of the burden would fall on owners as a one-off wealth loss at announcement rather than as a permanent flow; conversely, abolishing council tax would be expected to capitalise into higher property prices. The empirical literature does not settle the question: \citet{hoj2018} find full capitalisation in Danish municipal data, while \citet{nielsson2024} estimate a precise zero on later Danish variation and attribute it to pass-through to tenants under rent regulation. Our capitalisation figures should therefore be read as the textbook upper bound on the wealth effect, not as a forecast. Section~\ref{sec:robustness} quantifies the first of these: full capitalisation at discount rates of 3--5 per cent implies a fall in land values of 15--26 per cent, or \pounds 1.15--1.92 trillion, concentrated on the top wealth decile. Our flow-based accounting is best interpreted as the impact incidence in the initial years of the reform, and it understates how sharply the reform bears on land wealth.

\paragraph{Regional land shares.} Equation~(\ref{eq:land}) applies a single land share $\lambda_r$ to every property in a region. In practice land shares vary within regions --- they are higher in urban centres and lower for new-build dwellings on inexpensive land --- so we are likely to understate within-region heterogeneity in LVT liabilities, although decile-level averages should be less affected.

\paragraph{Corporate land allocation.} We allocate \pounds 2.06 trillion of corporate land to households in proportion to their measured corporate wealth. This abstracts from foreign ownership of UK corporate equity, which would shift part of the burden abroad and render the reform more favourable to UK households than we report, and it assumes full pass-through of the corporate-sector tax to shareholders rather than to workers or consumers.

\paragraph{Northern Ireland.} Northern Ireland levies domestic rates rather than council tax; the central specification leaves rates unchanged, so Northern Irish households pay LVT in addition to their existing property tax. Section~\ref{sec:robustness} reports the Great Britain-only alternative, which raises the required rate to 0.79 per cent and leaves the distributional pattern intact. A comprehensive UK-wide reform would fold domestic rates into the replacement, lowering the required rate slightly below our central figure.

\paragraph{Measurement.} Land values are imputed rather than observed: Wealth and Assets Survey property wealth is self-reported, the upper tail of the survey remains thin even after enhancement, and the National Balance Sheet's land estimates rest on the ONS's residual valuation method \citep{ons2025}. Council tax liabilities are modelled from regional band averages rather than parcel-level records.

\subsection{Policy implications and design options}

The most salient losing group comprises asset-rich, income-poor households --- typically older owner-occupiers of valuable homes, concentrated in London and the South East --- for whom an annual charge of 0.77 per cent of land value may represent a substantial share of cash income. Established remedies address this liquidity problem: deferral of the liability, with interest, until sale or death, as operated for property taxes in several US states and recommended by \citet{mirrlees2011} and \citet{muellbauer2018}; a phased transition over an extended period; or an exempt band of land value per household. Each option sacrifices some revenue or efficiency, but the underlying difficulty is one of liquidity rather than solvency, since the households concerned are by construction wealthy. The appendix quantifies the first option: households containing someone of state pension age owe 35.1 per cent of the total yield, \pounds 20.3 billion, so a universal pensioner deferral would postpone roughly a third of first-year receipts. That is the scale of the transitional financing requirement --- large, bounded, and a cash-flow cost rather than a revenue loss, since deferred liabilities are recovered with interest from estates.

An alternative reform family would replace council tax with a proportional tax on total property value (structures plus land), as proposed by \citet{muellbauer2005} and modelled by \citet{corlett2018}. A proportional property tax shares the LVT's correction of council tax's regressivity in value and obsolete valuations, and is administratively simpler because market prices are observed; however, it retains the disincentive to improve property and taxes the elastic (structures) margin alongside the inelastic one. We simulate this comparator directly in Section~\ref{sec:robustness}. At the 0.73 per cent rate that raises the same \pounds 57.6 billion, a proportional property tax produces fewer winners (64.4 against 68.2 per cent) and a shallower wealth gradient: the top wealth decile loses \pounds 4{,}217 rather than \pounds 5{,}752. The land base is more concentrated than the property base because location value is more concentrated than structure value, so for a given revenue target the LVT asks more of the largest landholders and less of everyone else. That is the quantitative content of the land/structures distinction, and it is the price paid for the proportional property tax's administrative advantage of observable market prices.

\subsection{Conclusion}

This paper has presented, to our knowledge, the first open-source household-level microsimulation of a revenue-neutral replacement of council tax with a land value tax for the United Kingdom. At the budget-neutral rate of 0.77 per cent, our results suggest that the reform would benefit roughly two-thirds of households, reduce after-housing-cost poverty modestly, and reallocate the burden of recurrent property taxation from council tax's regressive banded structure to the distribution of land wealth, with the largest losses concentrated in the top wealth decile. Measured income inequality and before-housing-cost poverty would be essentially unaffected --- a result that underscores the importance of examining incidence along the wealth as well as the income distribution, and a caution against reading the near-zero movement in headline statistics as evidence that little has changed. The principal obstacles to implementation are practical rather than analytical: valuing land separately from structures at scale, and managing the transition for asset-rich, income-poor owners. Both problems are long-standing, and both admit established policy responses. Extending the analysis to incorporate capitalisation, behavioural responses and parcel-level land valuation remains a subject for future research.
